\chapter*{Resumo}

\noindent\textbf{Autoria de atividades pedonais com estações sensíveis à capacidade e boids híbridos}

Muitas perguntas sobre edifícios e espaços públicos dependem do que as pessoas precisam, como uma secretária ou uma mesa, e do que acontece quando outra pessoa precisa dele ao mesmo tempo. Uma simulação que lhes responda tem de mostrar quando cada atividade começa de facto e porque se atrasa. Esta dissertação propõe uma estrutura em que horários em XML indicam quem faz o quê e quando, estações com capacidade limitada, fila e tempo de serviço decidem quando cada atividade ocorre, e uma camada de movimento híbrida segue caminhos de navegação, que os boids apenas desviam ligeiramente. Uma hora definida é, assim, um pedido, cujo início real depende do que a pessoa ainda faz, da distância e da disponibilidade do recurso, e um registo de regras, visitas e caminhadas explica cada atraso.

Implementámos a estrutura em C++ na Unreal Engine 5.4, com uma única implementação de pessoa numa rotina urbana diária com 50 pessoas, em dois centros comerciais, numa evacuação e em corredores de fluxo dirigido, com apenas uma ação nova em todos os cenários. Com a marcha e o relógio sincronizados, a avaliação mediu a origem da diferença entre os tempos definidos e os realizados. A espera por estações partilhadas dominou essa diferença nos centros comerciais, onde as filas foram 76\% do atraso no início das regras de estação e 92\% numa hora de almoço com um lugar por balcão, enquanto na rotina diária pesaram mais as caminhadas longas até destinos congestionados. As mudanças de capacidade comportaram-se como uma sequência de filas: em 30 execuções por configuração, um segundo ou terceiro lugar nas mesas reduziu a mediana da espera por visita de 0,84 para 0,51 minutos simulados, enquanto mais lugares nos balcões apenas deslocaram a fila para as mesas, como se viu na hora de almoço, com esperas de vários minutos. Estas tendências mantiveram-se com dois passos, mas não a sua dimensão. Os boids espalharam os grupos em dois de três corredores e aliviaram um estrangulamento, e um só computador anima 50 a 100 pessoas com uma média acima de 30 imagens por segundo, embora as mais lentas fiquem abaixo desse valor. Sem validação com peões reais, os resultados descrevem o modelo, não edifícios reais.

\paragraph{Palavras-chave} Autoria de regras, Boids, Recursos partilhados, Simulação de multidões, Unreal Engine

\cleardoublepage

\chapter*{Abstract}

\noindent\textbf{Authoring Pedestrian Activities with Capacity-Aware Stations and Hybrid Boids}

Many questions about buildings and public spaces are about what people need, such as a desk or a table, and what happens when someone else needs it at the same time. A simulation that answers them must show when each activity really starts and why it starts later than planned. This thesis presents a rule-authoring framework for pedestrian activities, in which XML schedules state who does what and from when, stations with a limited capacity, a first-come, first-served queue and a service time decide when each activity can take place, and a hybrid movement layer follows navigation paths while the boids forces only turn the heading slightly. An authored time is therefore a request, whose actual start depends on what the person is still doing, how far they must walk and whether the resource is free, and a recorder logs every rule, visit and walk to explain each delay.

We built the framework in C++ in Unreal Engine 5.4 and used one implementation of a person for a daily urban routine with 50 people, two shopping center layouts, an emergency evacuation and a set of directed-flow corridors, with only one new action across all scenarios. With walking and the clock advancing together, the evaluation measured where the gap between authored and realized times comes from. Waiting for shared stations was the largest part of that gap in the shopping centers, where queueing made up 76\% of the start delay of station rules and 92\% in a lunch rush with one place per counter, while in the daily routine long walks to crowded target points took a larger share. Capacity changes behaved as a pipeline of queues: over 30 runs per setting, a second or third place at the food tables lowered the median wait per station visit from 0.84 to 0.51 simulated minutes, while extra places at the food counters only moved the queue to the tables, and the lunch rush showed the same shift with waits of several minutes. These directions held at two frame steps, although their sizes changed. The boids blend spread the cohorts in two of three test corridors and eased a bottleneck, and one desktop computer animates 50 to 100 people at a mean above 30 frames per second, although its slowest frames fall below that. Without validation against real pedestrians, the results describe the model rather than real buildings.

\paragraph{Keywords} Boids, crowd simulation, rule authoring, shared resources, Unreal Engine

\cleardoublepage
